\chapter{Institutions and Historical Records}
\index{institutions}\index{archives}

\section{What Kind of Object Is an Institution?}

An institution is a compression architecture for collective history. Courts compress legal history into precedent and doctrine. Governments compress political history into laws and regulations. Archives compress documentary history into retrievable records. The trajectory space---sequences of decisions, records, laws, and organizational transitions---is substantially more tractable than the experiential trajectory space of Chapter~11.

\section{Institutional Trajectory Spaces}

Several types arise: legal (cases, statutes, judicial interpretations), political (elections, legislative acts), organizational (hiring, strategy, culture), and scientific (experimental results, publications, citations). Each supports different counterfactual questions.

\textbf{Diagnostic Question 1: Mostly affirmative.} Precedent depth for legal systems mirrors branch depth from Chapter~6, providing a tractable interaction order. Legal citation networks \cite{LandesPosner1976} provide an empirical proxy.

\section{Stability: Diagnostic Question 2}

Institutional influence exhibits heavy-tailed decay. Constitutional provisions from centuries ago retain binding authority; Roman law concepts persist across modern jurisdictions. Scientific citation half-lives \cite{Price1965} are finite, suggesting polynomial decay. Religious and cultural institutions may exhibit stability failure under temporal interaction order.

\textbf{Diagnostic Question 2: Conditional, polynomial stability likely.}

\section{Institutional Compressions: Diagnostic Question 3}

Institutional compressions are explicit and inspectable: legal archives, accounting records, scientific publications, oral traditions. Sufficiency is high for contemporary institutional queries; conditional for historical counterfactual queries not anticipated during compression design.

\section{Failure Modes and Institutions}

\begin{itemize}
  \item \textbf{Archive destruction} $\leftrightarrow$ sufficiency failure (irreversible)
  \item \textbf{Historical revisionism} $\leftrightarrow$ Type I degeneracy (compression collapses distinct trajectories)
  \item \textbf{Bureaucratic overgeneralization} $\leftrightarrow$ structural misalignment
  \item \textbf{Institutional lock-in} $\leftrightarrow$ stability failure for reform queries
  \item \textbf{Lost oral tradition} $\leftrightarrow$ compression-operator destruction
\end{itemize}

\section{Diagnostic Audit and Open Obligations}

\begin{center}\small
\begin{tabular}{lll}
\toprule
Question & Status & Notes \\
\midrule
Decomposition exact? & Mostly affirmative & Precedent depth tractable \\
Stability summable? & Conditional & Polynomial likely; varies by domain \\
Compression sufficient? & Affirmative (designed queries) & Conditional (historical queries) \\
Depths aligned? & Variable & Legal systems better than others \\
Degeneracy avoided? & Active concern & Revisionism demonstrates Type I \\
\bottomrule
\end{tabular}
\end{center}

Open obligations: (1) measure institutional influence decay; (2) characterize precedential equivalence; (3) measure archival sufficiency for counterfactual queries; (4) develop alignment-based archival design principles; (5) identify boundary between polynomial and non-summable stability for cultural institutions \cite{David1985,Arthur1994}.
